%% ========================================================================
%% Chapter 1: Pengenalan Sistem Operasi dan Instalasi Linux
%% ========================================================================

\chapter{Pengenalan Sistem Operasi dan Instalasi Linux}
\label{ch:introduction}

\chapterhero{blue}{Pahami fondasi sistem operasi, pilihan distribusi Linux, proses instalasi, dan verifikasi awal agar sistem siap dipakai sebagai lingkungan belajar.}{\faDesktop\quad operating system}{installation}{verification}

Setiap kali komputer dinyalakan, ada satu program yang bekerja bahkan sebelum aplikasi
lain dapat terbuka. Program tersebut mengatur seluruh perangkat keras, membagi waktu
prosesor di antara banyak tugas, dan memastikan setiap aplikasi mendapat ruang memori
yang aman. Program itu adalah \textbf{sistem operasi}. Tanpa sistem operasi, setiap
program harus menulis sendiri kode untuk berbicara dengan kartu jaringan, layar, dan
disk. Pendekatan itu tidak efisien dan hampir mustahil diterapkan secara luas.

Buku ini menggunakan Linux sebagai media pembelajaran. Linux menjalankan lebih dari 90
persen server di internet, seluruh 500 superkomputer terkencang di dunia, dan sebagian
besar infrastruktur cloud global. Memahami Linux berarti memahami fondasi teknologi
modern.

\textbf{Yang Akan Kamu Pelajari}

Setelah menyelesaikan bab ini, kamu akan mampu:

\begin{enumerate}
    \item menjelaskan definisi dan fungsi utama sistem operasi;
    \item membedakan jenis-jenis sistem operasi dan arsitektur kernel beserta contohnya;
    \item menelusuri garis besar evolusi sistem operasi dari era awal hingga saat ini;
    \item menginstalasi Ubuntu Server 24.04 LTS di \term{virtual machine} VirtualBox;
    \item menjelaskan konsep partisi disk dan tipe \term{filesystem} di Linux;
    \item melakukan konfigurasi awal dan memverifikasi hasil instalasi.
\end{enumerate}

\section{Definisi dan Fungsi Sistem Operasi}
\label{sec:os-basics}

Sistem operasi adalah perangkat lunak yang bertindak sebagai perantara antara pengguna
dan perangkat keras komputer. Bayangkan sebuah meja kerja. Di atas meja terdapat dokumen
yang sedang digarap, alat tulis, dan ruang untuk bergerak. Sistem operasi adalah manajer
meja kerja itu: ia menentukan dokumen mana yang boleh ada di atas meja sekarang, di mana
alat tulis disimpan, dan siapa yang boleh menyentuh benda mana.

Sistem operasi memiliki lima fungsi utama:

\begin{description}[noitemsep,topsep=2pt,parsep=0pt]
    \item[Manajemen Proses] Sistem operasi mengatur semua program yang berjalan. Ia
    menentukan giliran program mendapat waktu prosesor, membuat proses baru saat
    dibutuhkan, dan mengakhiri proses saat selesai. Di Linux, perintah \cmd{ps} dan
    \cmd{top} digunakan untuk melihat daftar proses yang sedang berjalan.

    \item[Manajemen Memori] Setiap program yang aktif membutuhkan ruang di RAM. Sistem
    operasi membagi-bagikan ruang itu, memastikan satu program tidak menyerobot area
    milik program lain, dan memindahkan data ke disk saat RAM mulai penuh.

    \item[Manajemen Berkas] Sistem operasi menyediakan struktur direktori sehingga data
    dapat disimpan, diberi nama, dan ditemukan kembali. Ia juga mengatur izin akses:
    siapa boleh membaca, mengubah, atau menjalankan sebuah berkas.

    \item[Manajemen I/O] Perangkat seperti keyboard, printer, dan kartu jaringan
    bekerja dengan cara yang sangat berbeda satu sama lain. Sistem operasi menyediakan
    lapisan abstraksi melalui \term{driver} sehingga program tidak perlu memahami cara kerja
    internal setiap perangkat.

    \item[Keamanan dan Proteksi] Sistem operasi memverifikasi identitas pengguna saat
    login, membatasi akses ke sumber daya berdasarkan izin yang ditetapkan, dan mencatat
    aktivitas penting untuk keperluan audit keamanan.
\end{description}

\begin{notebox}
Tanpa sistem operasi, setiap program harus dilengkapi \term{driver} sendiri untuk setiap
perangkat keras yang ada. Bayangkan harus menulis ulang kode pengendali layar setiap
kali membuat aplikasi baru. Sistem operasi menghilangkan pekerjaan berulang ini.
\end{notebox}

\section{Jenis Sistem Operasi}
\label{sec:os-categories}

Sistem operasi dikembangkan untuk berbagai jenis perangkat dengan kebutuhan yang berbeda.
Tabel~\ref{tab:os-categories} merangkum kategori utama yang beredar saat ini.

\begin{table}[tbp]
    \centering
    \caption{Kategori sistem operasi berdasarkan platform}
    \label{tab:os-categories}
    \begin{tblr}{
        colspec = {lX[l]l},
        width = \linewidth,
        hlines,
        vlines,
        row{1} = {font=\bfseries},
        cells = {font=\small}
    }
        Kategori & Karakteristik & Contoh \\
        Desktop OS &
            Antarmuka grafis untuk komputer pribadi &
            Windows 11, macOS, Ubuntu Desktop \\
        Server OS &
            Melayani banyak klien; biasanya tanpa antarmuka grafis &
            Ubuntu Server, RHEL, Windows Server \\
        Mobile OS &
            Layar sentuh dan efisiensi daya baterai &
            Android, iOS \\
        Real-Time OS &
            Respons dalam batas waktu yang sangat ketat &
            FreeRTOS, QNX, VxWorks \\
        Embedded OS &
            Sumber daya terbatas; untuk IoT dan router &
            Embedded Linux, OpenWRT \\
    \end{tblr}
\end{table}

Buku ini berfokus pada kategori \textbf{Server OS}, khususnya Linux. Keterampilan yang
diperoleh juga berlaku langsung di Desktop OS berbasis Linux dan Embedded OS yang
menggunakan kernel Linux yang sama.

\section{Arsitektur Kernel}
\label{sec:kernel-architecture}

Kernel adalah inti dari sistem operasi. Program ini selalu berada di memori sejak
komputer dinyalakan hingga dimatikan, dan memiliki akses langsung ke semua perangkat
keras. Semua bagian lain dari sistem operasi, termasuk antarmuka pengguna dan utilitas,
berada di luar kernel dan harus meminta layanan kepada kernel untuk mengakses perangkat
keras.

\subsection{Kernel Space and User Space}

Sistem operasi modern membagi memori menjadi dua wilayah yang terpisah secara ketat:
Pemisahan hubungan antarwilayah ini diringkas pada \cref{fig:os-layers}.
Istilah \textit{ring} berasal dari tingkat hak akses prosesor pada arsitektur CPU.
Semakin kecil nomornya, semakin tinggi hak aksesnya. Pada praktik sistem operasi modern,
dua tingkat yang paling penting adalah \textbf{Ring 0} untuk kernel dan \textbf{Ring 3}
untuk aplikasi biasa.

\begin{description}[noitemsep,topsep=2pt,parsep=0pt]
    \item[\term{Kernel space} (Ring 0)] Wilayah dengan hak akses tertinggi. Hanya kode kernel
    yang berjalan di sini. Kode di \term{kernel space} dapat membaca dan menulis ke semua
    alamat memori, serta berbicara langsung ke perangkat keras.

    \item[\term{User space} (Ring 3)] Wilayah tempat aplikasi biasa berjalan. Akses ke
    perangkat keras dibatasi sepenuhnya. Aplikasi harus melewati mekanisme yang disebut
    \textbf{system call} untuk meminta layanan dari kernel.
\end{description}

\begin{figure}[tbp]
    \centering
    \begin{tikzpicture}[node distance=1.2cm, scale=0.95]
        \node[draw, rectangle, minimum width=11cm, minimum height=1.4cm,
              fill=blue!10, align=center] (user) {
            \textbf{User Space (Ring 3)}\\[0.2em]
            \small Browser, Editor Teks, Media Player, Terminal
        };
        \node[draw, rectangle, minimum width=11cm, minimum height=0.8cm,
              below=1cm of user, fill=yellow!20] (syscall) {
            \textbf{Antarmuka System Call}
        };
        \node[draw, rectangle, minimum width=11cm, minimum height=1.8cm,
              below=1cm of syscall, fill=red!10, align=center] (kernel) {
            \textbf{Kernel Space (Ring 0)}\\[0.2em]
            \small Penjadwal Proses, Manajemen Memori, Sistem Berkas, Driver Perangkat
        };
        \node[draw, rectangle, minimum width=11cm, minimum height=1.2cm,
              below=1cm of kernel, fill=gray!20, align=center] (hw) {
            \textbf{Perangkat Keras}\\[0.2em]
            \small CPU, RAM, Disk, Kartu Jaringan, GPU
        };
        \draw[->, thick] (user.south) -- (syscall.north)
            node[midway, right, font=\small] {System Call};
        \draw[->, thick] (syscall.south) -- (kernel.north);
        \draw[->, thick] (kernel.south) -- (hw.north)
            node[midway, right, font=\small] {Akses Langsung};
    \end{tikzpicture}
    \caption{Lapisan arsitektur sistem operasi}
    \label{fig:os-layers}
\end{figure}

System call adalah mekanisme resmi satu-satunya yang digunakan aplikasi untuk meminta
layanan kernel. Saat sebuah program membuka berkas, misalnya, ia memanggil system call
\texttt{open()}. Kernel kemudian memvalidasi izin, berinteraksi dengan \term{driver} disk, dan
mengembalikan hasilnya kepada program yang memanggil.

\subsection{Jenis Arsitektur Kernel}

Ada tiga pendekatan utama dalam merancang kernel sistem operasi. Perbandingan
karakteristik utamanya dirangkum pada \cref{tab:kernel-types}, sedangkan gambaran
letak komponen layanannya ditunjukkan pada \cref{fig:kernel-comparison}.

\begin{table}[tbp]
    \centering
    \caption{Perbandingan jenis arsitektur kernel}
    \label{tab:kernel-types}
    \begin{tblr}{
        colspec = {lX[l]X[l]X[l]},
        width = \linewidth,
        hlines,
        vlines,
        row{1} = {font=\bfseries},
        cells = {font=\small}
    }
        Aspek & Monolitik & Mikrokernel & Hibrida \\
        Deskripsi &
            Semua layanan berjalan di \term{kernel space} dalam satu \term{binary} besar &
            Hanya layanan esensial di kernel; sisanya berjalan di \term{user space} &
            Menggabungkan pendekatan monolitik dan mikrokernel \\
        Kinerja & Sangat tinggi & Lebih rendah karena overhead komunikasi & Tinggi \\
        Keandalan &
            Satu bug dapat merusak seluruh sistem &
            Kegagalan layanan tidak merusak kernel &
            Baik \\
        Contoh & Linux, BSD & QNX, MINIX & Windows NT, macOS \\
    \end{tblr}
\end{table}

Linux menggunakan arsitektur monolitik dengan dukungan modul yang dapat dimuat secara
dinamis. Pendekatan ini memberi Linux kinerja tinggi sekaligus fleksibilitas karena
\term{driver} dapat ditambahkan tanpa harus mengompilasi ulang seluruh kernel.

\begin{figure}[tbp]
    \centering
    \resizebox{0.78\textwidth}{!}{%
    \begin{tikzpicture}
        % Monolitik
        \node[draw, rectangle, minimum width=3.5cm, minimum height=5cm]
            (mono) at (0,0) {};
        \node[above] at (mono.north) {\textbf{Monolitik}};
        \node[draw, fill=red!20, minimum width=3cm, minimum height=4.5cm]
            at (mono.center) {};
        \node[text width=2.5cm, align=center] at (mono.center)
            {Kernel\\[0.5em]Semua Layanan};
        \node[below=0.1cm of mono.south, text width=3.5cm, align=center,
            font=\small] {Linux, BSD};
        % Mikrokernel
        \node[draw, rectangle, minimum width=3.5cm, minimum height=5cm]
            (micro) at (5,0) {};
        \node[above] at (micro.north) {\textbf{Mikrokernel}};
        \node[draw, fill=red!20, minimum width=3cm, minimum height=1cm]
            at (micro.center) {\small Kernel Minimal};
        \node[draw, fill=blue!20, minimum width=3cm, minimum height=1.2cm]
            at ([yshift=1.5cm]micro.center) {\small Layanan};
        \node[below=0.1cm of micro.south, text width=3.5cm, align=center,
            font=\small] {MINIX, QNX};
        % Hibrida
        \node[draw, rectangle, minimum width=3.5cm, minimum height=5cm]
            (hybrid) at (10,0) {};
        \node[above] at (hybrid.north) {\textbf{Hibrida}};
        \node[draw, fill=red!20, minimum width=3cm, minimum height=2.5cm]
            at ([yshift=-0.5cm]hybrid.center) {\small Kernel + Beberapa};
        \node[draw, fill=blue!20, minimum width=3cm, minimum height=1.2cm]
            at ([yshift=1.3cm]hybrid.center) {\small Layanan};
        \node[below=0.1cm of hybrid.south, text width=3.5cm, align=center,
            font=\small] {Windows, macOS};
    \end{tikzpicture}%
    }
    \caption{Perbandingan arsitektur kernel: blok inti menunjukkan layanan di
    \term{kernel space}, sedangkan blok terpisah di bagian atas menunjukkan layanan di
    \term{user space}}
    \label{fig:kernel-comparison}
\end{figure}

\section{Evolusi Sistem Operasi}
\label{sec:os-history}

Sistem operasi berkembang seiring kemajuan perangkat keras dan kebutuhan pengguna.
Memahami garis besar sejarahnya menjelaskan mengapa sistem operasi modern dirancang
seperti yang ada sekarang. Tonggak utamanya diringkas pada \cref{tab:os-history}.

\begin{table}[tbp]
    \centering
    \caption{Garis besar evolusi sistem operasi}
    \label{tab:os-history}
    \begin{tblr}{
        colspec = {lX[l]},
        width = \linewidth,
        hlines,
        vlines,
        row{1} = {font=\bfseries},
        cells = {font=\small}
    }
        Era & Perkembangan Utama \\
        1950-an: Pemrosesan Awal &
            Komputer berukuran besar tanpa OS; program dijalankan satu per satu
            menggunakan kartu plong (punch card), tanpa interaksi pengguna langsung \\
        1969: Lahirnya Unix &
            Ken Thompson dan Dennis Ritchie di Bell Labs menciptakan Unix: sederhana,
            portabel, dan ditulis dalam bahasa C dengan filosofi \textit{do one thing well} \\
        1980-an: Era PC &
            MS-DOS (1981) untuk IBM PC; Apple Macintosh (1984) memperkenalkan antarmuka
            grafis; Windows 1.0 (1985) hadir sebagai lapisan di atas DOS \\
        1991: Linux &
            Linus Torvalds dari Finlandia merilis kernel Linux 0.01 sebagai perangkat
            lunak bebas dan terbuka; kompatibel dengan Unix \\
        1993: Windows NT &
            Microsoft merilis Windows NT dengan arsitektur kernel hibrida khusus untuk
            pasar korporat dan enterprise \\
        2007--2008: Era Mobile &
            iOS (2007) dan Android (2008) mendefinisikan ulang komputasi dengan
            pendekatan layar sentuh dan mobile-first \\
        2010-an: Era Cloud &
            Linux mendominasi cloud computing; Docker (2013) merevolusi cara aplikasi
            dikemas dan didistribusikan \\
    \end{tblr}
\end{table}

Unix menjadi fondasi hampir semua sistem operasi modern. Linux diturunkan dari filosofi
Unix, macOS dibangun di atas kernel berbasis Unix, dan Android menggunakan kernel Linux.
Perdebatan terkenal tahun 1992 antara Andrew Tanenbaum (pencipta MINIX, pendukung
mikrokernel) dan Linus Torvalds (pendukung desain monolitik) memperlihatkan bahwa kedua
pendekatan memiliki tempatnya masing-masing. Linux mendominasi server dengan desain
monolitik-modular, sementara QNX dengan mikrokernel mendominasi sistem otomotif.

\section{Mengapa Linux?}
\label{sec:why-linux}

Linux dipilih sebagai platform utama dalam buku ini karena beberapa alasan yang bersifat
praktis dan relevan dengan industri:

\begin{enumerate}
    \item \textbf{Gratis dan terbuka.} Kode sumber Linux tersedia untuk siapa saja.
    Tidak ada biaya lisensi. Setiap bagian sistem dapat dipelajari langsung dari
    kodenya sendiri.

    \item \textbf{Dominasi di server dan cloud.} Lebih dari 90 persen server web di
    internet menjalankan Linux. Platform cloud besar seperti AWS, Google Cloud, dan Azure
    menggunakan Linux sebagai sistem operasi utama mereka.

    \item \textbf{Konsep yang dapat dipindahkan.} Keterampilan Linux berlaku di semua
    sistem berbasis Unix, termasuk macOS dan BSD. Menguasai Linux berarti memiliki
    fondasi yang kuat untuk mempelajari sistem operasi lain.

    \item \textbf{Transparansi.} Linux tidak menyembunyikan cara kerjanya. Berkas
    konfigurasi dapat dibaca secara langsung, log sistem tersedia, dan setiap proses
    dapat diamati dengan perintah standar.

    \item \textbf{Nilai karier.} Permintaan tenaga kerja dengan keahlian Linux terus
    meningkat, terutama di bidang DevOps, administrasi sistem, dan rekayasa cloud.
\end{enumerate}

Buku ini menggunakan \textbf{Ubuntu Server 24.04 LTS} sebagai distribusi Linux yang
dipilih. Ubuntu Server dikenal stabil, memiliki dokumentasi yang lengkap, dan mendukung
pembaruan keamanan selama lima tahun. Ubuntu juga merupakan distribusi yang paling banyak
digunakan di lingkungan cloud komersial.

\begin{notebox}
Linux bukan satu sistem operasi tunggal. Linux adalah nama kernel. Berbagai distribusi
seperti Ubuntu, Debian, Fedora, dan Arch Linux semuanya menggunakan kernel Linux yang
sama, tetapi berbeda dalam \term{package manager}, konfigurasi bawaan, dan
target penggunanya.
\end{notebox}

\section{Instalasi Ubuntu Server di VirtualBox}
\label{sec:installation}

Praktek pada bab ini menggunakan \textbf{\term{virtual machine}}. \term{Virtual machine}
adalah komputer yang disimulasikan di dalam komputer fisik menggunakan perangkat
lunak. Pendekatan ini aman untuk belajar: jika terjadi kesalahan konfigurasi, \term{virtual machine}
dapat dikembalikan ke kondisi semula tanpa memengaruhi sistem utama sama sekali.

\begin{notebox}
Seluruh praktek pada bab ini dijalankan dari antarmuka teks (command line) Ubuntu
Server. Tidak ada antarmuka grafis. Ini adalah kondisi normal untuk server produksi.
\end{notebox}

\textbf{Prasyarat}: Sebelum memulai instalasi, pastikan hal-hal berikut sudah tersedia:

\begin{itemize}
    \item VirtualBox versi 6.1 atau lebih baru, terinstalasi di komputer utama
    \item Berkas ISO Ubuntu Server 24.04 LTS, diunduh dari situs resmi Ubuntu
    \item Komputer utama dengan RAM minimal 4 GB dan ruang disk kosong 30 GB
\end{itemize}

\begin{tipbox}
Ubuntu LTS (Long Term Support) berarti Canonical memberikan pembaruan keamanan dan
perbaikan bug selama lima tahun sejak rilis. Untuk server, selalu gunakan versi LTS
agar sistem tetap stabil dalam jangka panjang.
\end{tipbox}

\subsection*{Praktek 1.1 Membuat Mesin Virtual}

Buat sebuah \term{virtual machine} baru di VirtualBox untuk menginstalasi Ubuntu Server.

\begin{enumerate}
    \item Buka VirtualBox, lalu klik tombol \textbf{New} di bagian atas jendela.
    \item Isi kolom konfigurasi:
    \begin{itemize}
        \item Name: \texttt{Ubuntu-Server-Lab}
        \item Type: \texttt{Linux}
        \item Version: \texttt{Ubuntu (64-bit)}
    \end{itemize}
    \item Tetapkan memori (RAM) sebesar \textbf{2048 MB} (2 GB).
    \item Buat hard disk virtual baru dengan ukuran \textbf{25 GB}, tipe VDI, dan
    alokasi dinamis.
    \item Setelah \term{virtual machine} terbentuk, buka \textbf{Settings} lalu navigasi ke
    \textbf{Storage}.
    \item Pada Controller: IDE, tambahkan optical drive dan pilih berkas ISO Ubuntu
    Server yang sudah diunduh.
    \item Pastikan tab \textbf{Network} menggunakan mode \textbf{NAT} agar \term{virtual machine}
    dapat mengakses internet.
    \item Klik \textbf{Start} untuk menyalakan \term{virtual machine}.

    \textit{Amati:} Mesin virtual melakukan boot dari ISO Ubuntu Server dan menampilkan
    menu instalasi. Proses ini identik dengan menyalakan komputer fisik dari media
    instalasi bootable.
\end{enumerate}

\subsection*{Praktek 1.2 Proses Instalasi Ubuntu Server}

Ikuti langkah-langkah berikut setelah menu instalasi terbuka:

\begin{enumerate}
    \item Pilih bahasa: \textbf{English}.
    \item Pilih tata letak keyboard yang sesuai dengan keyboard yang digunakan.
    \item Pada layar konfigurasi jaringan, biarkan pengaturan DHCP otomatis berlaku.
    \item Lewati pengaturan proxy (biarkan kolom kosong).
    \item Pada bagian partisi, pilih \textbf{Use an entire disk} lalu konfirmasi untuk
    melanjutkan. Pengaturan ini akan membuat partisi secara otomatis.
    \item Isi profil pengguna seperti contoh berikut:
    \begin{lstlisting}[caption={Contoh isian profil pengguna}]
Your name: Budi Santoso
Server name: ubuntu-server
Username: budi
Password: [gunakan kata sandi yang kuat]
    \end{lstlisting}
    \item Aktifkan pilihan \textbf{Install OpenSSH server} agar sistem dapat diakses
    dari jarak jauh melalui terminal.
    \item Lewati pilihan snap tambahan dan tunggu proses instalasi selesai.
    \item Saat muncul pesan \textit{Installation complete}, pilih \textbf{Reboot Now}.

    \textit{Amati:} Proses instalasi membutuhkan waktu sekitar 15 hingga 20 menit,
    tergantung kecepatan koneksi internet. VirtualBox secara otomatis melepas ISO dari
    drive virtual saat sistem melakukan restart.
\end{enumerate}

\subsection*{Praktek 1.3 Login Pertama dan Verifikasi Sistem}

Setelah sistem restart, lakukan login pertama dan verifikasi bahwa instalasi berhasil.

\begin{enumerate}
    \item Masukkan nama pengguna dan kata sandi yang dibuat saat instalasi.
    \item Jalankan perintah berikut untuk memverifikasi informasi sistem:

    \begin{lstlisting}[language=bash, caption={Verifikasi informasi sistem}]
# Tampilkan versi kernel yang berjalan
uname -r

# Tampilkan informasi distribusi Linux
cat /etc/os-release

# Tampilkan alamat IP yang diperoleh dari DHCP
ip addr show
    \end{lstlisting}

    \textit{Amati:} Perintah \cmd{uname -r} menampilkan nomor versi kernel Linux yang
    sedang aktif. Perintah \cmd{cat /etc/os-release} menampilkan informasi distribusi.
    Pastikan nama distribusi menunjukkan Ubuntu 24.04 LTS.
\end{enumerate}

\begin{tipbox}
Setelah instalasi berhasil, segera buat snapshot di VirtualBox. Matikan \term{virtual machine},
buka menu \textbf{Snapshots}, klik \textbf{Take}, lalu beri nama \textit{Fresh Install}.
Snapshot ini memungkinkan pemulihan ke kondisi instalasi bersih kapan pun diperlukan
tanpa harus menginstalasi ulang dari awal.
\end{tipbox}

\section{Konfigurasi Awal Sistem}
\label{sec:post-installation}

Setelah sistem berhasil diinstalasi, ada beberapa langkah konfigurasi awal yang harus
dilakukan sebelum sistem siap digunakan untuk praktek selanjutnya.

\subsection*{Praktek 1.4 Memperbarui Sistem}

Pembaruan sistem memastikan semua paket yang terinstalasi berada pada versi terbaru
dengan tambalan keamanan yang tersedia.

\begin{lstlisting}[language=bash, caption={Memperbarui paket sistem}]
# Unduh informasi paket terbaru dari repository
sudo apt update

# Perbarui semua paket ke versi terbaru
sudo apt upgrade -y

# Hapus paket yang tidak lagi dibutuhkan
sudo apt autoremove -y
\end{lstlisting}

\textit{Amati:} Perintah \cmd{apt update} tidak menginstalasi apa pun. Perintah ini
hanya memperbarui daftar paket yang tersedia dari server \term{repository}. Perintah
\cmd{apt upgrade} baru melakukan instalasi berdasarkan daftar yang sudah diperbarui.

\subsection*{Praktek 1.5 Instalasi Utilitas Dasar}

Instalasi beberapa utilitas yang akan sering digunakan sepanjang praktek di buku ini:

\begin{lstlisting}[language=bash, caption={Instalasi utilitas dasar}]
sudo apt install -y \
    curl \
    wget \
    vim \
    htop \
    tree \
    net-tools \
    neofetch
\end{lstlisting}

Setelah instalasi selesai, jalankan \cmd{neofetch} untuk melihat ringkasan informasi
sistem:

\begin{lstlisting}[language=bash, caption={Menampilkan informasi sistem}]
neofetch
\end{lstlisting}

\textit{Amati:} Neofetch menampilkan logo distribusi Linux beserta informasi versi OS,
kernel, CPU, RAM, dan disk dalam satu tampilan ringkas. Perhatikan kolom memori: berapa
persen RAM yang sudah terpakai oleh sistem dasar Ubuntu Server tanpa aplikasi tambahan.

\subsection*{Praktek 1.6 Verifikasi Koneksi dan Sumber Daya Sistem}

\begin{lstlisting}[language=bash, caption={Pemeriksaan jaringan dan sumber daya}]
# Uji konektivitas internet dengan mengirim 4 paket ke DNS Google
ping -c 4 8.8.8.8

# Tampilkan penggunaan disk dalam format mudah dibaca
df -h

# Tampilkan penggunaan memori dan swap
free -h

# Tampilkan proses yang sedang berjalan secara interaktif
htop
\end{lstlisting}

\textit{Amati:} Perintah \cmd{ping} mengirimkan empat paket data ke alamat 8.8.8.8
(server DNS Google). Jika respons kembali dengan nilai \textit{time} yang terukur,
koneksi internet berfungsi dengan baik. Perintah \cmd{df -h} menampilkan penggunaan
setiap partisi dalam satuan GB atau MB yang mudah dibaca. Tekan \keystroke{q} untuk
keluar dari \cmd{htop}.

\section{Partisi Disk dan Filesystem}
\label{sec:disk-partitioning}

Sebuah disk dapat dibagi menjadi beberapa bagian yang disebut \textbf{partisi}. Bayangkan
sebuah lemari arsip dengan laci-laci terpisah: satu laci untuk dokumen kerja, satu untuk
dokumen pribadi, dan satu untuk arsip lama. Partisi bekerja dengan cara yang sama. Setiap
partisi bersifat independen, dapat menggunakan \term{filesystem} yang berbeda, dan dapat
diproteksi secara terpisah.

\subsection{Tabel Partisi: MBR dan GPT}

Ada dua standar tabel partisi yang digunakan saat ini. Perbedaannya diringkas pada
\cref{tab:mbr-gpt}.

\begin{table}[tbp]
    \centering
    \caption{Perbandingan MBR dan GPT}
    \label{tab:mbr-gpt}
    \begin{tblr}{
        colspec = {lX[l]X[l]},
        width = \linewidth,
        hlines,
        vlines,
        row{1} = {font=\bfseries},
        cells = {font=\small}
    }
        Aspek & MBR (Legacy) & GPT (Modern) \\
        Firmware & BIOS & UEFI \\
        Batas ukuran disk & 2 TB & 9,4 ZB (praktis tidak terbatas) \\
        Jumlah partisi & Maksimal 4 partisi primer & Hingga 128 partisi \\
        Penggunaan & Sistem lama & Semua sistem produksi modern \\
    \end{tblr}
\end{table}

Instalasi Ubuntu Server 24.04 LTS pada sistem modern menggunakan GPT secara otomatis.

\subsection{Linux Partition Layout}

Saat memilih \textit{Use an entire disk} pada Praktek 1.2, Ubuntu membuat
\term{partition layout} berikut secara otomatis:

\begin{description}[noitemsep,topsep=2pt,parsep=0pt]
    \item[\texttt{/boot/efi}] Partisi EFI berukuran kecil (512 MB) untuk menyimpan
    bootloader GRUB yang digunakan firmware UEFI saat komputer dinyalakan.
    \item[\texttt{/}] Partisi root yang menampung seluruh sistem operasi, program, dan
    konfigurasi. Ini adalah partisi utama Linux.
    \item[swap] Area disk yang digunakan sebagai perluasan RAM virtual. Saat RAM penuh,
    data yang tidak aktif dipindahkan ke swap untuk memberi ruang bagi proses yang aktif.
\end{description}

\subsection{Tipe Filesystem}

Setiap partisi diformat dengan \term{filesystem} tertentu yang menentukan bagaimana data
disimpan dan diorganisasikan. Contoh yang paling umum pada instalasi Linux diringkas pada
\cref{tab:filesystems}.

\begin{table}[tbp]
    \centering
    \caption{Tipe filesystem umum di Linux}
    \label{tab:filesystems}
    \begin{tblr}{
        colspec = {lX[l]l},
        width = \linewidth,
        hlines,
        vlines,
        row{1} = {font=\bfseries},
        cells = {font=\small}
    }
        Filesystem & Karakteristik & Digunakan untuk \\
        ext4 & Stabil, matang, mendukung journaling & Partisi root (\texttt{/}) dan data \\
        swap & Area virtual memory di disk & Partisi swap \\
        FAT32 / vfat & Kompatibel dengan semua firmware & Partisi \texttt{/boot/efi} \\
    \end{tblr}
\end{table}

\subsection*{Praktek 1.7 Memeriksa Partisi dan Filesystem}

Setelah sistem terinstalasi, periksa \term{partition layout} yang dibuat secara otomatis.

\begin{lstlisting}[language=bash, caption={Memeriksa partisi dan filesystem}]
# Tampilkan disk layout dan partisi dalam bentuk pohon
lsblk

# Tampilkan detail partisi termasuk ukuran dan tipe
sudo fdisk -l

# Tampilkan tipe filesystem setiap partisi
lsblk -f
\end{lstlisting}

\textit{Amati:} Perintah \cmd{lsblk} menampilkan semua disk dan partisi dalam struktur
pohon. Perhatikan nama perangkat seperti \texttt{sda1}, \texttt{sda2}, dan seterusnya,
serta ukuran masing-masing. Perintah \cmd{lsblk -f} menambahkan kolom FSTYPE yang
menunjukkan \term{filesystem} setiap partisi, misalnya \texttt{ext4} untuk partisi root
dan \texttt{vfat} untuk partisi EFI.

\section{Rangkuman}
\label{sec:rangkuman}

Dalam bab ini, kamu telah mempelajari:

\begin{itemize}[noitemsep,topsep=2pt]
    \item \textbf{Fungsi dan jenis sistem operasi}: sistem operasi mengelola proses,
    memori, berkas, I/O, dan keamanan; terdapat lima kategori OS berdasarkan platform,
    mulai dari desktop dan server hingga mobile, real-time, dan embedded.
    \item \textbf{Arsitektur kernel}: kernel berjalan di \term{kernel space} dengan hak akses
    penuh; aplikasi di \term{user space} berkomunikasi melalui system call; tiga arsitektur
    utama adalah monolitik (Linux), mikrokernel (QNX), dan hibrida (Windows, macOS).
    \item \textbf{Evolusi sistem operasi}: Unix (1969) menjadi fondasi sistem operasi
    modern; Linux (1991) diturunkan dari filosofi Unix dan kini mendominasi server serta
    infrastruktur cloud global.
    \item \textbf{Instalasi Ubuntu Server}: \term{virtual machine} di VirtualBox menyediakan
    lingkungan pembelajaran yang aman; instalasi mencakup pembuatan VM, pemasangan OS,
    dan verifikasi sistem melalui perintah \cmd{uname -r} dan \cmd{ip addr show}.
    \item \textbf{Partisi dan filesystem}: disk dibagi menjadi partisi-partisi
    independen; GPT menggantikan MBR pada sistem modern; Linux menggunakan ext4 untuk
    partisi root, FAT32 untuk EFI, dan area swap sebagai \term{virtual memory}.
    \item \textbf{Konfigurasi awal}: pembaruan sistem dengan \cmd{apt update} dan
    \cmd{apt upgrade}, instalasi utilitas dasar, serta verifikasi partisi dengan
    \cmd{lsblk} memastikan sistem siap untuk praktek selanjutnya.
\end{itemize}

\section{Latihan}
\label{sec:tugas}

\textbf{Instruksi Umum}: Kerjakan seluruh latihan secara mandiri. Catat langkah
penting, simpan tangkapan layar bila diperlukan, lalu rangkum hasilnya sebagai
dokumentasi pribadi.

\subsubsection*{Latihan 1.1 Instalasi dan Verifikasi Sistem}
Lakukan instalasi Ubuntu Server 24.04 LTS di VirtualBox sesuai langkah pada
Praktek 1.1 hingga 1.3. Setelah instalasi selesai:
\begin{enumerate}
    \item Jalankan \cmd{uname -r} dan catat versi kernel yang ditampilkan.
    \item Jalankan \cmd{cat /etc/os-release} dan catat nama serta versi distribusi.
    \item Jalankan \cmd{ip addr show} dan catat alamat IP yang diperoleh.
    \item Ambil screenshot dari output setiap perintah di atas.
\end{enumerate}


\subsubsection*{Latihan 1.2 Konfigurasi Awal dan Verifikasi}
Setelah sistem terinstalasi, lakukan langkah-langkah konfigurasi awal berikut:
\begin{enumerate}
    \item Perbarui sistem dengan \cmd{sudo apt update} diikuti \cmd{sudo apt upgrade -y}.
    \item Instal paket \cmd{neofetch} dan jalankan. Ambil screenshot hasilnya.
    \item Jalankan \cmd{df -h} dan \cmd{free -h}. Catat kapasitas disk dan memori yang
    tersedia pada sistem.
    \item Buat snapshot di VirtualBox dengan nama \textit{Fresh Install} setelah semua
    langkah selesai.
\end{enumerate}


\subsubsection*{Latihan 1.3 Analisis Konsep Sistem Operasi}
Jawab pertanyaan berikut berdasarkan materi yang telah dipelajari pada bab ini:
\begin{enumerate}
    \item Sebuah perusahaan membutuhkan sistem operasi untuk tiga keperluan: server web,
    perangkat medis dengan respons waktu ketat, dan laptop karyawan. Kategori OS apa
    yang paling sesuai untuk masing-masing keperluan tersebut? Jelaskan alasannya.
    \item Jelaskan perbedaan antara \term{kernel space} dan \term{user space}. Mengapa aplikasi
    tidak diizinkan mengakses perangkat keras secara langsung?
    \item Unix lahir pada 1969 dan Linux pada 1991. Jelaskan hubungan antara keduanya
    dan mengapa filosofi Unix masih relevan hingga era cloud computing saat ini.
    \item Jalankan perintah \cmd{lsblk -f} pada sistem Ubuntu Server yang sudah
    terinstalasi. Catat nama partisi, tipe \term{filesystem}, dan ukuran setiap partisi.
    Identifikasi partisi mana yang merupakan root (\texttt{/}), partisi EFI, dan area
    swap. Sertakan output perintah sebagai screenshot.
\end{enumerate}


\section{Referensi}
\label{sec:referensi}

\begin{itemize}
    \item Silberschatz, A., Galvin, P. B., Gagne, G. (2018). \textit{Operating System
    Concepts} (10th ed.). Wiley.

    \item Tanenbaum, A. S., Bos, H. (2014). \textit{Modern Operating Systems} (4th
    ed.). Pearson.

    \item Ward, B. (2021). \textit{How Linux Works} (3rd ed.). No Starch Press.

    \item Canonical Ltd. (2024). \textit{Ubuntu Server Documentation}.
    \url{https://ubuntu.com/server/docs}

    \item The Linux Kernel Organization. (2024). \textit{The Linux Kernel Archives}.
    \url{https://www.kernel.org}
\end{itemize}
